% v11 RQ1/RQ2 on student-authored repositories (WP H1 selection, WP H2 evaluation). Numbers: numbers-wph2 macros.
\subsection{Student-Authored Repositories}\label{sec:student}
The real-history projects of Section~\ref{sec:real} are professional or our own. As a stand-in for student code we selected \hTwoRepos{} public Python repositories that describe themselves as capstone, senior-design, final-year or course projects, by a frozen search and screening protocol: team authorship, at least 2,000 lines of Python and 60 first-parent commits, and at least 90\% of commits before December 2022. They are student work, but they are not supervised resubmissions, and commits are not submissions. The original commit-pair selection and syntax-tree reference were applied unchanged, yielding \hTwoN{} observations from \hTwoPairs{} commit pairs (\hTwoTargetedN{} edit-targeted) after \hTwoExcluded{} undecidable targets were excluded. No threshold, guard or design was chosen on these repositories, so every method is evaluated here on held-out data (Table~\ref{tab:student}).

% Generated by tools/v11/build_wph2_tex.py; do not edit.
\begin{table}[t]
\caption{Student-Authored Repositories (Held Out)}\label{tab:student}
\centering\footnotesize\setlength{\tabcolsep}{3pt}
\begin{tabular}{@{}lrrrrrr@{}}\toprule
 & \multicolumn{4}{c}{All observations} & \multicolumn{2}{c}{Edit-targeted} \\ \cmidrule(lr){2-5}\cmidrule(l){6-7}
Method & Acc & Missed & Silent & Extra & Acc & Missed \\ \midrule
Deployed cascade & 91.5 & 3.1 & 1.5 & 0.0 & 73.6 & 2.6 \\
Pre-existing-copy guard & 91.5 & 0.2 & 0.4 & 4.4 & 73.6 & 0.1 \\
All guards, $\theta=.85$ & 92.8 & 0.2 & 0.4 & 4.7 & 78.2 & 0.1 \\
Cascade, $\theta=.55$ & 94.6 & 1.2 & 1.1 & 0.2 & 84.1 & 0.9 \\
All guards, $\theta=.55$ & 95.9 & 0.1 & 0.4 & 4.7 & 88.7 & 0.0 \\
Cascade v1 & 91.3 & 8.2 & 2.6 & 0.0 & 73.0 & 7.8 \\
Contextual tracker & 94.9 & 0.5 & 0.2 & 3.3 & 85.4 & 0.3 \\
Diff drift & 76.1 & 0.1 & 0.1 & 5.4 & 20.9 & 0.0 \\
Exact quote & 77.8 & 0.5 & 0.2 & 3.2 & 25.7 & 0.3 \\
Line number & 43.6 & 0.1 & 0.0 & 47.3 & 13.1 & 0.1 \\
\bottomrule
\end{tabular}
\tablenote{}{Percentages over 7,882 observations (1,857 edit-targeted) from 20 repositories. Acc: file/start, including correct absence. Missed: changed or gone targets without a re-check request. Silent: wrong placement without a request. Extra: requests on preserved targets.}
\end{table}


The deployed cascade placed \hTwoDepAcc\% of the targets correctly (95\% commit-pair cluster interval \hTwoDepAccLo--\hTwoDepAccHi\%), and the contextual tracker \hTwoTrkAcc\%. The tracker's lead (paired cluster interval of deployed minus tracker: \hTwoDiffLo{} to \hTwoDiffHi{} points) held in \hTwoTrkLeadRepos{} of the \hTwoRepos{} repositories, with \hTwoTieRepos{} ties, and widened on edit-targeted observations (\hTwoDepTgtAcc\% against \hTwoTrkTgtAcc\%). This replicates the real-history finding of Section~\ref{sec:realtracker} on independent code: the cascade's location recovery is not state of the art. Here the tracker was given Git's rename detection. The common-input comparison of Section~\ref{sec:locator} gives every locator the anchor's original path only; under that rule the tracker placed \wiStuTrQAcc\% of the same targets, the cascade's figure, because \wiStuTrRenamed{} of its placements were on renamed files. Its lead on edit-targeted observations remains (\wiStuTgtTrQAcc\%).

The safety results replicate as well, and here they are held out. The deployed cascade failed to request \hTwoDepMiss\% of the needed re-checks (interval \hTwoDepMissLo--\hTwoDepMissHi\%) and placed \hTwoDepSilent\% of the issues on wrong code without a request. The pre-existing-copy guard alone reduced missed re-checks to \hTwoAGMiss\%; all three guards reduced them to \hTwoAllMiss\% and silent misplacement to \hTwoAllSilent\%, while raising accuracy to \hTwoAllAcc\%. Their cost was \hTwoAllFalse\% extra requests on preserved targets, concentrated on edit-targeted ones (\hTwoAllTgtFalse\%). With the threshold of $0.55$ chosen by leave-one-project-out calibration on the original projects, all three guards reached \hTwoAllCalAcc\% accuracy with \hTwoAllCalMiss\% missed re-checks, close to the tracker's \hTwoTrkAcc\% and \hTwoTrkMiss\%. The guards were designed from the original histories' failures; that they transfer to unseen student repositories is the strongest evidence in this paper that they address general mechanisms rather than the cases that motivated them.
